\section{Introduction}

Après avoir présenté le cadre du projet, ce chapitre détaille les besoins de PharmacieConnect. L'objectif
est d'identifier les acteurs, de préciser les fonctionnalités attendues et de fixer les principales règles
de gestion. Cette étape sert de base à la conception technique présentée dans le chapitre suivant.

\section{Identification des acteurs}

Le système distingue plusieurs acteurs. Chacun intervient avec un rôle différent et un niveau d'accès
adapté à ses besoins.

\begin{longtable}{L{4cm}L{10cm}}
    \caption{Acteurs de la plateforme PharmacieConnect}\\
    \toprule
    \textbf{Acteur} & \textbf{Rôle dans le système} \\
    \midrule
    \endfirsthead
    \toprule
    \textbf{Acteur} & \textbf{Rôle dans le système} \\
    \midrule
    \endhead
    Administrateur & Gère les pharmacies, les plannings de garde, les médicaments, les imports CSV, les indicateurs, les utilisateurs autorisés et les journaux d'audit. \\
    Assistant ou utilisateur autorisé & Intervient sur un périmètre limité, notamment pour mettre à jour les pharmacies et les gardes selon les droits accordés. \\
    Utilisateur mobile ou citoyen & Recherche une pharmacie, consulte les pharmacies de garde, utilise la géolocalisation, consulte les médicaments et gère ses préférences. \\
    Système & Valide les données, applique les règles métier, sécurise les accès, calcule les distances, journalise les actions et expose les API. \\
    \bottomrule
\end{longtable}

\section{Besoins fonctionnels}

Les besoins fonctionnels décrivent les services que la plateforme doit offrir. Ils sont organisés par
module afin de mieux séparer les responsabilités.

\begin{longtable}{L{4cm}L{10cm}}
    \caption{Besoins fonctionnels principaux}\\
    \toprule
    \textbf{Module} & \textbf{Fonctionnalités attendues} \\
    \midrule
    \endfirsthead
    \toprule
    \textbf{Module} & \textbf{Fonctionnalités attendues} \\
    \midrule
    \endhead
    Authentification & Connexion, inscription, vérification email, rafraîchissement de session, déconnexion et consultation du profil. \\
    Recherche de pharmacies & Recherche par texte, consultation de la liste des pharmacies, détail d'une pharmacie et recherche par proximité. \\
    Pharmacies de garde & Consultation des gardes par date, import des plannings, création, modification et suppression côté administration. \\
    Gestion des pharmacies & Création, modification, suppression, consultation, pagination et import CSV des pharmacies. \\
    Gestion des médicaments & Import CSV, recherche par nom commercial, DCI ou code PCT, consultation du détail et pagination. \\
    Administration & Tableau de bord, gestion des assistants, contrôle des autorisations, carte administrative, analytics et audit logs. \\
    Application mobile & Accueil, carte, calendrier des gardes, médicaments, paramètres, langues, thème sombre, favoris et notifications. \\
    \bottomrule
\end{longtable}

\section{Besoins non fonctionnels}

Les besoins non fonctionnels concernent la qualité attendue du système. Ils ne décrivent pas une page ou
une action précise, mais ils influencent directement la conception et la réalisation.

\begin{itemize}
    \item \textbf{Sécurité} : protéger les routes administratives, hacher les mots de passe, utiliser des jetons JWT,
    vérifier les rôles et tracer les actions sensibles.
    \item \textbf{Performance} : paginer les réponses, limiter les tailles d'import, éviter les réponses trop
    volumineuses et utiliser le cache lorsque la configuration le permet.
    \item \textbf{Ergonomie} : proposer une application mobile simple, une administration claire et une navigation
    adaptée aux opérations fréquentes.
    \item \textbf{Maintenabilité} : séparer les routes, services, modèles, schémas, composants et écrans afin de
    faciliter l'évolution du projet.
    \item \textbf{Disponibilité} : rendre les informations importantes accessibles à travers des endpoints publics
    et prévoir une configuration adaptée au déploiement.
    \item \textbf{Fiabilité des données} : valider les fichiers CSV, contrôler les champs obligatoires, signaler les
    erreurs et éviter les doublons lorsque des identifiants métier existent.
\end{itemize}

\section{Diagramme de cas d'utilisation}

Le diagramme de cas d'utilisation présente les interactions principales entre les acteurs et la plateforme.
L'image finale devra être ajoutée lorsque le diagramme UML sera disponible.

\placeholderfigure
{Diagramme de cas d'utilisation}
{Image UML à insérer : interaction entre administrateur, assistant, utilisateur mobile et système.}
{fig:use-case-placeholder}

\section{Description textuelle des cas d'utilisation importants}

\subsection{Cas d'utilisation : authentification}

\begin{longtable}{L{4cm}L{10cm}}
    \caption{Description textuelle du cas d'utilisation Authentification}\\
    \toprule
    \textbf{Élément} & \textbf{Description} \\
    \midrule
    \endfirsthead
    \toprule
    \textbf{Élément} & \textbf{Description} \\
    \midrule
    \endhead
    Acteurs & Administrateur, assistant, utilisateur mobile \\
    Objectif & Permettre à un utilisateur d'accéder aux fonctionnalités correspondant à son profil. \\
    Précondition & L'utilisateur possède un compte valide ou suit le processus d'inscription. \\
    Scénario nominal & L'utilisateur saisit ses identifiants, le backend vérifie les informations, génère les jetons JWT et retourne les données de session. \\
    Postcondition & L'utilisateur est authentifié et peut accéder aux fonctionnalités autorisées. \\
    Exceptions & Identifiants invalides, compte désactivé, email non vérifié ou jeton expiré. \\
    \bottomrule
\end{longtable}

\subsection{Cas d'utilisation : recherche d'une pharmacie}

\begin{longtable}{L{4cm}L{10cm}}
    \caption{Description textuelle du cas d'utilisation Recherche d'une pharmacie}\\
    \toprule
    \textbf{Élément} & \textbf{Description} \\
    \midrule
    \endfirsthead
    \toprule
    \textbf{Élément} & \textbf{Description} \\
    \midrule
    \endhead
    Acteur & Utilisateur mobile ou citoyen \\
    Objectif & Trouver une pharmacie à partir d'un mot-clé, d'un gouvernorat ou de la position GPS. \\
    Précondition & L'application mobile est installée et l'API est accessible. \\
    Scénario nominal & L'utilisateur saisit une recherche ou autorise la géolocalisation, l'application interroge l'API, le backend filtre les pharmacies et retourne les résultats. \\
    Postcondition & Les pharmacies correspondantes sont affichées sous forme de liste ou sur la carte. \\
    Exceptions & Position indisponible, recherche vide, absence de résultat ou indisponibilité temporaire de l'API. \\
    \bottomrule
\end{longtable}

\subsection{Cas d'utilisation : consultation des pharmacies de garde}

\begin{longtable}{L{4cm}L{10cm}}
    \caption{Description textuelle du cas d'utilisation Consultation des gardes}\\
    \toprule
    \textbf{Élément} & \textbf{Description} \\
    \midrule
    \endfirsthead
    \toprule
    \textbf{Élément} & \textbf{Description} \\
    \midrule
    \endhead
    Acteur & Utilisateur mobile ou citoyen \\
    Objectif & Consulter les pharmacies de garde pour une date donnée. \\
    Précondition & Des plannings de garde sont enregistrés dans la base. \\
    Scénario nominal & L'utilisateur choisit une date, l'application interroge l'API des gardes et affiche les pharmacies disponibles. \\
    Postcondition & Les gardes correspondant à la date choisie sont affichées. \\
    Exceptions & Date invalide, absence de planning ou erreur de communication avec l'API. \\
    \bottomrule
\end{longtable}

\subsection{Cas d'utilisation : gestion des pharmacies}

\begin{longtable}{L{4cm}L{10cm}}
    \caption{Description textuelle du cas d'utilisation Gestion des pharmacies}\\
    \toprule
    \textbf{Élément} & \textbf{Description} \\
    \midrule
    \endfirsthead
    \toprule
    \textbf{Élément} & \textbf{Description} \\
    \midrule
    \endhead
    Acteurs & Administrateur, assistant autorisé \\
    Objectif & Ajouter, modifier, consulter ou supprimer les données des pharmacies. \\
    Précondition & L'acteur est authentifié et possède les droits nécessaires. \\
    Scénario nominal & L'acteur ouvre la page des pharmacies, effectue l'opération souhaitée, le backend valide les données et enregistre la modification. \\
    Postcondition & Les données sont mises à jour et l'action est journalisée si nécessaire. \\
    Exceptions & Données invalides, pharmacie introuvable, droits insuffisants ou périmètre régional non autorisé. \\
    \bottomrule
\end{longtable}

\subsection{Cas d'utilisation : import CSV}

\begin{longtable}{L{4cm}L{10cm}}
    \caption{Description textuelle du cas d'utilisation Import CSV}\\
    \toprule
    \textbf{Élément} & \textbf{Description} \\
    \midrule
    \endfirsthead
    \toprule
    \textbf{Élément} & \textbf{Description} \\
    \midrule
    \endhead
    Acteurs & Administrateur, assistant autorisé selon le périmètre \\
    Objectif & Importer un ensemble de pharmacies, de gardes ou de médicaments à partir d'un fichier CSV. \\
    Précondition & L'acteur est authentifié et le fichier respecte le format attendu. \\
    Scénario nominal & L'acteur charge le fichier, le backend vérifie l'extension, la taille, les colonnes, les valeurs et enregistre les lignes valides. \\
    Postcondition & Les données valides sont enregistrées et un bilan d'import est retourné. \\
    Exceptions & Colonnes manquantes, lignes invalides, fichier trop volumineux, doublons ou données hors périmètre régional. \\
    \bottomrule
\end{longtable}

\subsection{Cas d'utilisation : gestion des utilisateurs et assistants}

\begin{longtable}{L{4cm}L{10cm}}
    \caption{Description textuelle du cas d'utilisation Gestion des utilisateurs et assistants}\\
    \toprule
    \textbf{Élément} & \textbf{Description} \\
    \midrule
    \endfirsthead
    \toprule
    \textbf{Élément} & \textbf{Description} \\
    \midrule
    \endhead
    Acteur & Administrateur ou super administrateur \\
    Objectif & Créer, modifier, activer, désactiver ou limiter les comptes administratifs et assistants. \\
    Précondition & L'acteur dispose d'un rôle autorisé. \\
    Scénario nominal & L'administrateur renseigne les informations du compte, choisit le rôle et le périmètre éventuel, puis valide l'opération. \\
    Postcondition & Le compte est créé ou mis à jour avec les permissions correspondantes. \\
    Exceptions & Email déjà utilisé, rôle invalide, droits insuffisants ou informations manquantes. \\
    \bottomrule
\end{longtable}

\section{Règles de gestion}

Les règles de gestion retenues sont les suivantes :

\begin{itemize}
    \item une route d'administration ne peut être appelée que par un compte authentifié et autorisé ;
    \item un assistant ne peut intervenir que dans le périmètre régional qui lui est attribué ;
    \item les mots de passe doivent être hachés avant leur enregistrement ;
    \item les jetons d'accès expirés doivent être renouvelés à travers un jeton de rafraîchissement valide ;
    \item une pharmacie doit posséder au minimum un nom et des informations de localisation exploitables ;
    \item les coordonnées GPS doivent respecter les bornes valides de latitude et de longitude ;
    \item un planning de garde doit être associé à une date valide ;
    \item un fichier CSV importé doit respecter les colonnes obligatoires du module concerné ;
    \item les erreurs d'import doivent être retournées avec un message compréhensible ;
    \item les actions sensibles doivent être journalisées dans les logs d'audit.
\end{itemize}

\section{Conclusion}

Ce chapitre a précisé les acteurs, les besoins et les règles de gestion de PharmacieConnect. Ces éléments
permettent de passer à la conception de la solution en gardant un lien clair entre les fonctionnalités
attendues et les choix techniques.
